Tras clonar el repositorio, se copia el fichero de variables de entorno
de ejemplo y se generan las claves criptográficas:

\begin{verbatim}
git clone <url-del-repositorio> scdee
cd scdee
cp backend/.env.example backend/.env
make keygen
\end{verbatim}

El comando \texttt{make keygen} imprime por consola la clave secreta de
Django y la clave maestra de cifrado (\acs{aes}-256-\acs{gcm} para el
\acs{dni}). Ambos valores deben copiarse al fichero \texttt{backend/.env},
sustituyendo los marcadores de posición existentes.

Las variables del fichero \texttt{.env} se agrupan en seis bloques:
Django (\texttt{DJANGO\_SECRET\_KEY}, \texttt{DJANGO\_DEBUG},
\texttt{DJANGO\_ALLOWED\_HOSTS}), base de datos (\texttt{DB\_*}), Redis
y Celery (\texttt{REDIS\_URL}, \texttt{CELERY\_*}), MinIO
(\texttt{MINIO\_*}), correo (\texttt{EMAIL\_*}) y cifrado
(\texttt{ENCRYPTION\_MASTER\_KEY}). Las credenciales por defecto de la
base de datos son \texttt{scdee} / \texttt{scdee\_dev\_password}.
